\capitulo{3}{Conceptos teóricos}

En esta sección de la memoria, se pretenden recoger los conceptos teóricos más relevantes para la comprensión del proyecto.

\subsection{Aplicaciones web y arquitecturas cliente-servidor}

Una aplicación web es un sistema software al que se puede acceder a través de un navegador, más concretamente a través de una url, con lo que requiere tener conexión a Internet. Además está sustentada sobre una arquitectura cliente-servidor.

Esta arquitectura o modelo tiene funciona de la siguiente manera.
\begin{itemize}
	\item El cliente se encarga de la interacción con el usuario y la representación de la interfaz, es decir, es el que hace las peticiones de servicios al servidor y recibe las respuestas de estas.
	\item El servidor es le que centraliza la lógica de negocio, el acceso a los datos, y recibe y procesa las peticiones del cliente etnregándoles una respuesta.
\end{itemize}

La comunicación entre el cliente y el servidor se realiza por medio de .protocolos de red. En las aplicaciones web, el más común es \textit{\textbf{HTTP}} (\textit{Hyper Text Transfer Protocol}), un protocolo de comunicación que define cómo se intercambian las peticiones y respuestas entre cliente y servidor. Gracias a este mecanismo, el cliente puede solicitar o enviar recursos o información al servidor, mientras que este puede responder con datos o el resultado de una operación concreta.
\cite{clienteservidor:Fernan}

Esta separación de responsabilidades facilita el mantenimiento, la escalabilidad y la evolución de forma independiente de las distintas partes que conforman el sistema. En este proyecto se utiliza un framework distinto para el front (cliente) que para el back (servidor).

\subsection{API REST}

Una API (\textit{Application Programming Interface}) es un conjunto de reglas, protocolos y definiciones que permite la comunicación entre diferentes sistemas software. En el contexto de las aplicaciones web, una API permite que el \textit{frontend} y el \textit{backend} intercambien información de una forma estructurada, sin la necesidad de que ambos softwares estén implementados con la misma tecnologías.

En particular, una API REST (\textit{Representational State Transfer}) es un estilo de arquitectura ampliamente utilizado en el desarrollo de servicios web. Su objetivo principal es organizar la comunicación entre cliente y servidor mediante recursos identificados por URLs y operaciones estándar del protocolo HTTP. De esta forma permite que los recursos del sistema como, en el caso de este proyecto, un aula o una reserva, pueda ser modificado o consultado a través de una dirección concreta y utilizando unos métodos bien definidos.

Entre los métodos más comunes de HTTP en una API REST se encuentran:
\begin{itemize}
	\item \textit{GET}: utilizado para solicitar o recuperar información.
	\item \textit{POST}: envía datos al servidor para crear un nuevo recurso.
	\item \textit{PUT} o \textit{PATCH}: para modificaciones de recursos ya existentes.
	\item \textit{DELETE}: utilizado para la eliminación de recursos existentes.
\end{itemize}
Con esta organización las operaciones del sistema quedan estructuradas.
\cite{metodosHttp:Oscar}


\tablaSmall{Herramientas y tecnologías utilizadas en cada parte del proyecto}{l c c c c}{herramientasportipodeuso}
{ \multicolumn{1}{l}{Herramientas} & App AngularJS & API REST & BD & Memoria \\}{ 
	HTML5 & X & & &\\
	CSS3 & X & & &\\
	BOOTSTRAP & X & & &\\
	JavaScript & X & & &\\
	AngularJS & X & & &\\
	Bower & X & & &\\
	PHP & & X & &\\
	Karma + Jasmine & X & & &\\
	Slim framework & & X & &\\
	Idiorm & & X & &\\
	Composer & & X & &\\
	JSON & X & X & &\\
	PhpStorm & X & X & &\\
	MySQL & & & X &\\
	PhpMyAdmin & & & X &\\
	Git + BitBucket & X & X & X & X\\
	Mik\TeX{} & & & & X\\
	\TeX{}Maker & & & & X\\
	Astah & & & & X\\
	Balsamiq Mockups & X & & &\\
	VersionOne & X & X & X & X\\
} 